%!TEX root = ./ft-hd.tex
\subsection{Recovery at constant SNR}\label{sec:const-snr}
The algorithm is given by

\begin{algorithm}[H]
\caption{\textsc{RecoverAtConstantSNR}($\hat x, \chi, k, \e$)}\label{alg:const-snr} 
\begin{algorithmic}[1] 
\Procedure{RecoverAtConstantSNR}{$\hat x, \chi, k, \e$}
\State $B\gets  \GT\cdot k/(\e\alpha^d)$
\State Choose $\Sigma\in \gl$, $q\in \nsq$ uniformly at random, let $\pi:=(\Sigma, q)$ and let $H_r:=(\pi_r, B, F)$
\State Let $\A\gets $ $C\log\log N$ elements of $\nsq\times \nsq$ sampled uniformly at random with replacement
\State $\H\gets \{\mathbf{0}_d\}$, $\Delta\gets 2^{\lfloor \frac1{2}\log_2 \log_2 n\rfloor}$ \Comment{$\mathbf{0}_d$ is the zero vector in dimension $d$}
\For {$g=1$ to $\lceil \log_\Delta n \rceil$}
\State $\H\gets \H\cup \bigcup_{s=1}^d n \Delta^{-g} \cdot \mathbf{e}_s$ \Comment{$\mathbf{e}_s$ is the unit vector in direction $s$}
\EndFor
\For {$\h\in \H$}
\State $m(\wh{x}, H, a\star(\one, \h))\gets \Call{HashToBins}{\hat x, 0, (H, a\star (\one, \h))}$ for all $a\in \A, \h\in \H$
\EndFor
\State $L \gets \textsc{LocateSignal}\left(\chi^{(t)}, k, \{m(\wh{x}, H, a\star (\one, \h))\}_{a\in \A, \h\in \H}\right)$
\State $\chi' \gets \Call{EstimateValues}{\hat x,  \chi, 2k, \e, O(\log N), 0}$
\State $L'\gets $top $4k$ elements of $\chi'$
\State \textbf{return} $\chi+\chi'_{L'}$
\EndProcedure 
\end{algorithmic}
\end{algorithm}


Our analysis will use 
\begin{lemma}[Lemma~9.1 from~\cite{IKP}]\label{lm:comparison}
Let $x, z\in \C^n$ and $k\leq n$. Let $S$ contain the largest $k$ terms of $x$, and $T$ contain the largest $2k$ terms of $z$. Then $||x-z_T||_2^2\leq ||x-x_S||_2^2+4||(x-z)_{S\cup T}||_2^2$.
\end{lemma}


\noindent{\em {\bf Lemma~\ref{lm:const-snr}}
For any $\e>0$, $\hat x, \chi\in \C^N$, $x'=x-\chi$ and any integer $k\geq 1$ if $||x'_{[2k]}||_1\leq O(||x_{\nsq\setminus [k]}||_2\sqrt{k})$ and $||x'_{\nsq\setminus [2k]}||_2^2\leq ||x_{\nsq\setminus [k]}||_2^2$, the following conditions hold. If $||x||_\infty/\mu=N^{O(1)}$, then the output $\chi'$ of 
\Call{RecoverAtConstantSNR}{$\hat x, \chi, 2k, \e$} satisfies
$$
||x'-\chi'||^2_2\leq (1+O(\e))||x_{\nsq\setminus [k]}||_2^2
$$
with at least $99/100$ probability over its internal randomness. The sample complexity is $2^{O(d^2)}\frac1{\e}  k\log N$, and the runtime complexity is at most $2^{O(d^2)}\frac1{\e}  k \log^{d+1} N.$
}

\begin{remark}
We note that the error bound is in terms of the $k$-term approximation error of $x$ as opposed to the $2k$-term approximation error of $x'=x-\chi$.
\end{remark}
\begin{proof}
Let $S$ denote the top $2k$ coefficients of $x'$. We first derive bounds on the probability that an element $i\in S$ is not located.
Recall that by Lemma~\ref{lm:loc} for any $i\in S$ if
\begin{enumerate}
\item  $e^{head}_{i}(H, x', 0)<|x'_i|/20$;
\item  $e^{tail}_{i}(H, \A\star (\one, \h), x')< |x'_i|/20$ for all $\h\in \H$;
\item for every $s\in [1:d]$ the set $\A\star(\zero, \mathbf{e}_s)$ is balanced (as per Definition~\ref{def:balance}),
\end{enumerate}
then $i\in L$, i.e. $i$ is successfully located in \textsc{LocateSignal}. 
  

We now upper bound the probability that an element $i\in S$ is not located.  We let $\mu^2:=||x_{\nsq\setminus k}||_2^2/k$ to simplify notation.

\paragraph{Contribution from head elements.} We need to bound, for $i\in S$, the quantity
\begin{equation*}
\begin{split}
e^{head}_i(H, x', 0)&=G_{o_i(i)}^{-1}\cdot \sum_{j\in S\setminus \{i\}} G_{o_i(j)}|x'_j|.
\end{split}
\end{equation*}
Recall that $m(\wh{x}, H, a\star(\one, \h))=\Call{HashToBins}{\hat x, 0, (H, a\star (\one, \h))}$, and let $m:=m(\wh{x}, H, a\star(\one, \h))$ to simplify notation. By Lemma~\ref{lm:hashing}, {\bf (1)} one has 
\begin{equation}\label{eq:const-h}
\expect_H[\max_{a\in \nsq} \abs{G_{o_i(i)}^{-1}\omega^{-a^T\Sigma i}m_{h(i)} - (x'_{S})_i}]\leq \GS\cdot C^d ||x'_{S}||_1/B+\mu/N^2
\end{equation}
for a constant $C>0$. This yields
$$
\expect_H[e^{head}_i(H, x', 0)]\leq \GS \cdot C^d ||x'_{S}||_1/B\lesssim  \GS \cdot C^d \mu k/B\lesssim \alpha^d C^d \e \mu.
$$
by the choice of $B$ in \textsc{RecoverAtConstantSNR}. Now by Markov's inequality we have for each $i\in \nsq$ 
\begin{equation}\label{eq:const-eh-prob}
\prob_H[e^{head}_{i}(H, x', 0)>|x'_i|/20]\lesssim \alpha^d C^d \e \mu/|x'_i|\lesssim \alpha \e \mu/|x'_i|
\end{equation}
as long as $\alpha$ is smaller than a constant.

\paragraph{Contribution of tail elements}
We restate the definitions of $e^{tail}$ variables here for convenience of the reader (see~\eqref{eq:et-pi}, \eqref{eq:et-pi-a-h}, \eqref{eq:et-pi-a} and \eqref{eq:et}). 

We have
\begin{equation*}
e^{tail}_i(H, z, x):=\left|G_{o_i(i)}^{-1}\cdot \sum_{j\in \nsq\setminus S} G_{o_i(j)}x_j \omega^{z^T \Sigma (j-i)}\right|.
\end{equation*}

For any $\mathcal{Z}\subseteq \nsq$ we have
\begin{equation*}
e^{tail}_i(H, \mathcal{Z}, x):=\quant^{1/5}_{z\in \mathcal{Z}} \left|G_{o_i(i)}^{-1}\cdot \sum_{j\in \nsq\setminus S} G_{o_i(j)}x_j \omega^{z^T \Sigma (j-i)}\right|.
\end{equation*}
Note that the algorithm first selects sets $\A_r\subseteq \nsq\times \nsq$, and then accesses the signal at locations given by $\A_r\star (\one, \h), \h\in \H$ (after permuting input space).

The definition of $e^{tail}_i(H, \A_r, x')$ for permutation $\pi=(\Sigma, q)$ allows us to capture the amount of noise that our measurements for locating a specific set of bits of $\Sigma i$ suffer from. Since the algorithm requires all $\h\in \H$ to be not too noisy in order to succeed (see preconditions 2 and 3 of Lemma~\ref{lm:loc}), we have
\begin{equation*}
e^{tail}_i(H, \A, x')=40\mu_{H, i}(x)+\sum_{\h\in \H} \left|e^{tail}_i(H, \A\star (\one, \h), x')-40\mu_{H, i}(x')\right|_+
\end{equation*}
where for any $\eta\in \mathbb{R}$ one has $|\eta|_+=\eta$ if $\eta>0$ and $|\eta|_+=0$ otherwise. 


For each $i\in S$ we now define an error event $\E^*_i$ whose non-occurrence implies location of element $i$, and then show that for each $i\in S$ one has
\begin{equation}\label{eq:g}
\prob_{H, \A}[\E^*_i]\lesssim \frac{\alpha \e \mu^2}{|x'_i|^2}.
\end{equation}
Once we have ~\eqref{eq:g}, together with~\eqref{eq:const-eh-prob} it allows us to prove the main result of the lemma. In what follows we concentrate on proving 
~\eqref{eq:g}. Specifically, for each $i\in S$ define
\begin{equation*}
\E^*_i=\{(H, \A): \exists \h\in \H \text{~s.t.~} e^{tail}_i(H,\A\star (\one, \h), x')>|x'_i|/20\}.
\end{equation*}


Recall that $e^{tail}_i(H, z, x')=\Call{HashToBins}{\widehat{x_{\nsq\setminus S}}, \chi_{\nsq\setminus S}, (H, z)}$ by definition of the measurements $m$.  By Lemma~\ref{lm:hashing}, {\bf (3)} one has, for a uniformly random $z\in \nsq$, that 
$\expect_z[|e^{tail}_i(H, z, x')|^2|]=\mu_{H, i}^2(x')$. By Lemma~\ref{lm:hashing}, {\bf (2)}, we have that 
\begin{equation}\label{eq:expect-sigmaq}
\expect_H[\mu^2_{H, i}(x')] \leq  \GSS\cdot C^d \norm{2}{(x-\chi)_{\nsq\setminus S}}^2/B+\mu^2/N^2\leq  \alpha \e \mu^2.
\end{equation} 
Thus by Markov's inequality 
$$
\prob_{z}[e^{tail}_{i}(H, z, x')^2> (|x'_i|/20)^2]\leq \alpha\e (\mu_{H, i}(x'))^2/(|x'_i|/20)^2.
$$
Combining this with Lemma~\ref{lm:quant-exp}, we get for all $\tau\leq (1/20)(|x'_i|/20)$ and all $\h\in \H$
\begin{equation}\label{eq:xyz}
\prob_{\A}[\quant^{1/5}_{z\in \A\star (\one, \h)} e^{tail}_{i}(H, z, x')> |x'_i|/20| \mu_{H, i}^2(x')=\tau]<(4\tau/(|x'_i|/20))^{\Omega(|\A|)}.
\end{equation}


Equipped with the bounds above, we now bound $\prob[\E^*_i]$. To that effect, for each $\tau>0$ let the event $\E_i(\tau)$ be defined as 
$\E_i(\tau)=\{\mu_{H, i}(x')=\tau\}$. Note that since we assume that we operate on $O(\log n)$ bit integers, $\mu_{H, i}(x')$ takes on a finite number of values, and hence $\E_i(\tau)$ is well-defined. It is convenient to bound $\prob[\E_i^*]$ as a sum of three terms:
\begin{equation*}
\begin{split}
\prob_{H, \A}[\E^*_i]&\leq \prob_{H, \A}\left[\left.e^{tail}_{i}(H, \A, x')> |x'_i|/20\right|\bigcup_{\tau\leq \sqrt{\alpha \e}\mu}\E_i(\tau)\right] \\
&+\int_{\sqrt{\alpha \e}\mu}^{(1/8)(|x'_i|/20)} \prob_{H, \A}\left[\left.e^{tail}_{i}(H, \A, x')> |x'_i|/20\right.|\E_i(\tau)\right] \prob[\E_i(\tau)]d\tau\\
&+\int_{(1/8)(|x'_i|/20)}^\infty \prob[\E_i(\tau)]d\tau\\
\end{split}
\end{equation*}


We now bound each of the three terms separately for $i$ such that $|x'_i|/20\geq 2\sqrt{\alpha\e} \mu_{H, i}(x')$. This is sufficient for our purposes, as other elements only contribute a small amount of $\ell_2^2$ mass.
\begin{description}
\item[1.] By \eqref{eq:xyz} and a union bound over $\H$ the first term is bounded by 
\begin{equation}\label{eq:r-ub-1}
|\H|\cdot (\sqrt{\alpha \e}\mu/(|x'_i|/20))^{\Omega(|A|)}\leq \alpha \e \mu^2/|x'_i|^2\cdot |\H|\cdot 2^{-\Omega(|A|)}\leq \alpha\e \mu^2/|x'_i|^2.
\end{equation}
since $|\A|\geq C\log \log N$ for a sufficiently large constant $C>0$ in \textsc{RecoverAtConstantSNR}.

\item[2.] The second term, again by a union bound over $\H$ and using \eqref{eq:xyz}, is bounded by 
\begin{equation}\label{eq:const-snr-abc}
\begin{split}
&\int_{\sqrt{\alpha \e}\mu}^{(1/8)(|x'_i|/20)} |\H|\cdot (4\tau /(|x'_i|/20))^{\Omega(|A|)} \prob[\E_i(\tau)]d\tau\\
\leq &\int_{\sqrt{\alpha \e}\mu}^{(1/8)(|x'_i|/20)} |\H|\cdot (4\tau /(|x'_i|/20))^{\Omega(|A|)} (4\tau /(|x'_i|/20))^{2}\prob[\E_i(\tau)] d\tau\\
\end{split}
\end{equation}
Since $|\A|\geq C\log \log N$ for a sufficiently large constant $C>0$ and $(4\tau /(|x'_i|/20))\leq 1/2$ over the whole range of $\tau$ by our assumption that $|x'_i|/20\geq 2\sqrt{\alpha\e} \mu_{H, i}(x')$, we have
$$
|\H|\cdot (4\tau /(|x'_i|/20))^{\Omega(|A|)}\leq |\H|\cdot (1/2)^{\Omega(|A|)}=o(1)
$$ 
for each $\tau\in [\sqrt{\alpha \e}\mu, (1/8)(|x'_i|/20)]$. Thus, \eqref{eq:const-snr-abc} is upper bounded by
\begin{equation*}
\begin{split}
\int_{\sqrt{\alpha \e}\mu}^{(1/8)(|x'_i|/20)} (4\tau /(|x'_i|/20))^{2}\prob[\E_i(\tau)] d\tau\\
\lesssim \frac1{(|x'_i|/20)^{2}}\int_{\sqrt{\alpha \e}\mu}^{(1/8)(|x'_i|/20)} \tau^2\prob[\E_i(\tau)] d\tau\\
\leq \frac{\alpha \e \mu^2}{(|x'_i|/20)^2}\\
\end{split}
\end{equation*}
since 
$$
\int_{\sqrt{\alpha \e}\mu}^{(1/8)(|x'_i|/20)} \tau^2\prob[\E_i(\tau)] d\tau\leq \int_{0}^{\infty} \tau^2\prob[\E_i(\tau)] d\tau=\expect_H[\mu^2_{H, i}(x')]=O(\alpha)\e \mu^2
$$
by \eqref{eq:expect-sigmaq}.


\item [3.] For the third term we have
$$
\int_{(1/8)(|x'_i|/20)}^\infty \prob[\E_i(\tau)]d\tau=\prob[\mu_{H, i}(x')>(1/8)(|x'_i|/20)]\lesssim \frac{\alpha \e \mu^2}{|x'_i|^2}
$$
by Markov's inequality applied to \eqref{eq:expect-sigmaq}.
\end{description}
Putting the three estimates together, we get $\prob[\E_i^*]=\frac{O(\alpha)\e \mu^2}{|x'_i|^2}$. Together with \eqref{eq:const-eh-prob} this yields for $i\in S$
$$
\prob[i\not \in L]\lesssim \frac{\alpha \e \mu^2}{|x'_i|^2}+\frac{\alpha \e \mu}{|x'_i|}.
$$

In particular, 
\begin{equation*}
\begin{split}
\expect\left[\sum_{i\in S} |x'_i|^2\cdot \mathbf{1}_{i\in S\setminus L}\right]&\leq \sum_{i\in S} |x'_i|^2 \prob[i\not \in L]\\
&\lesssim \sum_{i\in S} |x'_i|^2 \left(\frac{\alpha \e \mu}{|x'_i|}+\frac{\alpha \e \mu^2}{|x'_i|^2}\right)\lesssim \alpha \e\mu^2 k,
\end{split}
\end{equation*}
where we used the assumption of the lemma that $||x'_{[2k]}||_1\leq O(||x_{\nsq\setminus [k]}||_2\sqrt{k})$ and $||x'_{\nsq\setminus [2k]}||_2^2\leq ||x_{\nsq\setminus [k]}||_2^2$ in the last line.
By Markov's inequality we thus have $\prob[||x'_{S\setminus L}||_2^2>\e \mu^2 k]<1/10$ as long as $\alpha$ is smaller than a constant. 

We now upper bound $||x'-\chi'||_2^2$. We apply Lemma~\ref{lm:comparison} to vectors $x'$ and $\chi'$ with sets $S$ and $L'$ respectively, getting 
\begin{equation*}
\begin{split}
||x'-\chi'_{L'}||_2^2&\leq ||x'-x'_S||_2^2+4||(x'-\chi')_{S\cup L'}||_2^2\\
&\leq ||x_{\nsq\setminus [k]}||_2^2+4||(x'-\chi')_{S\setminus L}||_2^2+4||(x'-\chi')_{S\cap L}||_2^2\\
&\leq ||x_{\nsq\setminus [k]}||_2^2+4\e \mu^2 k+4\e \mu^2 |S|\\
&\leq ||x_{\nsq\setminus [k]}||_2^2+O(\e \mu^2 k),\\
\end{split}
\end{equation*}
where we used the fact that $||(x'-\chi')_{S\cap L}||_\infty\leq \sqrt{\e} \mu$ with probability at least $1-1/N$ over the randomness used in \textsc{EstimateValues} by Lemma~\ref{lm:estimate-l1l2}, {\bf (3)}.
This completes the proof of correctness.


\paragraph{Sample complexity and runtime}The number of samples taken is bounded by $2^{O(d^2)}\frac1{\e} k\log N$ by Lemma~\ref{l:hashtobins}, the choice of $B$. The sampling complexity of the call to \textsc{EstimateValues} is at most $2^{O(d^2)}\frac1{\e} k\log N$. The runtime is bounded by $2^{O(d^2)}\frac1{\e}k \log^{d+1} N \log\log N$ for computing the measurements $m(\wh{x}, H, a\star (\one, \h))$ and $2^{O(d^2)}\frac1{\e} k \log^{d+1} N$ for estimation.
\end{proof}